\documentclass[tikz, border=10pt]{standalone}
\usepackage{tikz}
\usepackage{amsmath} % For \vec, \nabla, \partial commands
\renewcommand{\vec}[1]{\mathbf{#1}}
\usetikzlibrary{shapes.geometric, arrows.meta, fit}
\graphicspath{{../figure_primatives/}}

\begin{document}

\begin{tikzpicture}[
    block/.style={
        rectangle, draw, very thick,
        text width=4.5cm, align=center,
        minimum height=3cm, rounded corners=8pt,
        font=\large
    },
    small_block/.style={
        rectangle, draw, very thick,
        align=center, minimum height=1.2cm,
        font=\large,rounded corners=4pt
    },
    arrow/.style={-Latex, very thick}
]

% --- COORDINATE DEFINITIONS ---
\coordinate (c_func)      at (-8, 3);
\coordinate (c_phys_eq)   at (-8, -1);
\coordinate (c_geom)      at (-8, -5);
\coordinate (c_initial) at (-2,3);

\coordinate (c_pde_loss)  at (2, -1);
\coordinate (c_bc_loss)   at (2, -2.5);
\coordinate (c_data_loss) at (2, -4);
\coordinate (c_loss_label) at (2,0);

\coordinate (c_data)      at (2, -6);
\coordinate (c_colloc)    at (-3, -3);

\coordinate (c_optimizer) at (8, 0);
\coordinate (c_opt_func)  at (8, -4.5);
% ---------------------------------


% ---------------- NODES ----------------
\node[block] (func) at (c_func) {
    \textbf{Function Ansatz} \\
    $u(\vec{x}) = \sum_{i,j} c_{ij} \Phi_{ij}(\vec{x})$};

\node[block] (initial) at (c_initial) {
\textbf{Initial $\boldsymbol{c_{ij}}$}\\
\includegraphics[width=3cm]{schematic_b.png}
};

\node[block] (phys_eq) at (c_phys_eq) {
    \textbf{Physical Equation} \\
    $\nabla^2 u = 0 \quad \text{in } \Omega$ \\
    $u|_{\partial\Omega} = g$
};

\node[block] (geom) at (c_geom) {
    \textbf{Geometry} \\
    \includegraphics[width=3cm]{schematic_c.png}
};

% Middle: Loss Function components
\node[small_block] (pde_loss) at (c_pde_loss) {PDE Residuals};
\node[small_block] (bc_loss) at (c_bc_loss) {Boundary Residuals};
\node[small_block] (data_loss) at (c_data_loss) {Data Loss};
\node (loss_label) at (c_loss_label) {\large \textbf{Loss Function}};
\node[draw, very thick, rounded corners=8pt, inner sep=0.25cm,
      fit=(loss_label) (pde_loss) (bc_loss) (data_loss)] (loss) {};

% Inputs to the loss function
\node[small_block] (data) at (c_data) {\textbf{Data}};
\node[small_block] (colloc) at (c_colloc) {\textbf{Collocation Points}};

% Right side: Optimizer and Final Result
\node[block] (optimizer) at (c_optimizer) {
    \textbf{Optimizer} \\ (e.g., Adam, Levenberg-Marquardt)
};

\node[block] (opt_func) at (c_opt_func) {
    \textbf{Optimized $\boldsymbol{c_{ij}}$} \\
    \includegraphics[width=3cm]{schematic_a.png}
};


% ---------------- ARROWS ----------------

\coordinate (fork_top) at (-1.5, -.5);
\coordinate (fork_bottom) at (-1.5, -2);
% \draw (fork_top) -- (fork_bottom); % Draw the shared vertical line

% Draw paths from the sources (phys_eq, colloc) to this shared line.
\draw[arrow] (phys_eq.east) -- (-.2,-1);
\draw[very thick] (colloc.north) -- (-3,-1);

% \draw[arrow] (fork_top) -- (pde_loss.west);
% \draw[arrow] (fork_bottom) -- (bc_loss.west);

% Other arrows
\draw[arrow] (geom.east) -| (colloc.south); % A direct line is cleaner here
\draw[arrow] (data.north) -- (loss.south);
\draw[arrow] (loss.east) -- ++(.5,0) |- (optimizer.west);
\draw[arrow] (func.east) -- (initial.west);
\draw[arrow] (initial.east) -| (optimizer.north);
\draw[arrow] (optimizer.south) --(opt_func.north);

\end{tikzpicture}

\end{document}